% Capítulo 6 — Afinidades y coordenadas homogéneas
\chapter{Afinidades y coordenadas homogéneas}\label{cap:06}

\section{La idea}

Aflojamos la rigidez: una afinidad conserva rectas y paralelismo, y las coordenadas homogéneas permiten escribir la traslación como un producto de matrices. La formulación sigue lo dictado y los resúmenes del profesor; el paso a cada dimensión es desarrollo nuestro.

\subsection{Afinidades: conservar rectas y paralelismo}\label{sec:afinidades-primera}

{\small\vspace{1mm}\linkUnidad{04_afines_homogeneas}{Simulación interactiva: Unidad 6 — Afinidades y coordenadas homogéneas}\par}


Una afinidad permite deformaciones más generales. Un cuadrado puede convertirse en
paralelogramo y un círculo en elipse, pero las rectas siguen siendo rectas y las paralelas
siguen siendo paralelas cuando la transformación es invertible.

\begin{center}
\begin{tikzpicture}[scale=.86]
  \draw[teal,line width=1.1pt] (1,1) rectangle (3.7,3.7);
  \draw[rust,line width=1.1pt] (7,1)--(10.2,1.5)--(9.1,4.2)--(5.9,3.7)--cycle;
  \draw[-{Latex[length=2.8mm]},tealdark,line width=1pt]
    (4.3,2.4)--(5.4,2.4);
  \node[labelback,text=teal] at (2.35,.55) {cuadrado};
  \node[labelback,text=rust] at (8.05,.55) {paralelogramo};
\end{tikzpicture}
\end{center}

Los puntos fijos muestran dónde se ``ancla'' la transformación. Las direcciones propias
muestran líneas que conservan su dirección aunque puedan estirarse, contraerse o invertirse.
Las coordenadas homogéneas permiten tratar la traslación dentro del mismo lenguaje matricial.

\paragraph{Primera traducción algebraica.}
\[
  F(x)=Ax+t.
\]
La parte \(A\) deforma el paralelogramo de la base y \(t\) transporta su origen. Si \(A\) es
invertible, ninguna dirección completa colapsa en un punto.

\begin{idea}[En palabras del profesor]
``Una afinidad conserva paralelas las paralelas'' o ``una afinidad conserva rectas las
rectas y proporcionales las distancias sobre una recta'' (resumen de clase 3, cursada 2025).
\end{idea}

\subsection{Coordenadas homogéneas: puntos como rayos}\label{sec:homogeneas-primera}

{\small\vspace{1mm}\linkUnidad{04_afines_homogeneas}{Simulación interactiva: Unidad 6 — Afinidades y coordenadas homogéneas}\par}


En clase presentamos el plano afín como el plano \(w=1\) dentro de un espacio
vectorial tridimensional. El punto afín \((x,y)\) se representa mediante
\[
  (x,y)\longleftrightarrow
  \begin{pmatrix}x\\y\\1\end{pmatrix}.
\]
Sin embargo, en geometría proyectiva no importa el representante concreto, sino el rayo que
parte del origen:
\[
\begin{pmatrix}x\\y\\1\end{pmatrix}
\sim
\lambda
\begin{pmatrix}x\\y\\1\end{pmatrix},
\qquad \lambda\neq0.
\]
Todos esos vectores pertenecen al mismo subespacio vectorial de dimensión uno y cortan al
plano \(w=1\) en el mismo punto afín.

\begin{center}
\begin{tikzpicture}[scale=.92]
  \coordinate (o) at (1,1);
  \draw[axisline,-{Latex[length=2.4mm]}] (o)--(8.2,1) node[right] {$y$};
  \draw[axisline,-{Latex[length=2.4mm]}] (o)--(-.2,-.1) node[below left] {$x$};
  \draw[axisline,-{Latex[length=2.4mm]}] (o)--(1,6) node[above] {$w$};
  \draw[teallight,fill=teallight!55,line width=.8pt]
    (-.2,3.2)--(6.8,3.2)--(8.2,4.25)--(1.2,4.25)--cycle;
  \node[labelback,text=tealdark] at (7.1,4.55) {plano \(w=1\)};
  \coordinate (p) at (4.8,3.74);
  \fill[rust] (p) circle (2.7pt) node[labelback,above] {\((x,y,1)\)};
  \draw[rust,line width=1.05pt,-{Latex[length=2.5mm]}] (o)--($(o)!1.38!(p)$);
  \fill[ink] (o) circle (2.4pt) node[labelback,below] {origen};
  \node[labelback,font=\small,text=rust] at (3.4,2.1)
    {un rayo representa un punto};
\end{tikzpicture}
\end{center}

\paragraph{Puntos en el infinito.}
Una recta afín con dirección \(v\) puede escribirse
\[
  p(\mu)=p+\mu v.
\]
Homogéneamente:
\[
\begin{pmatrix}p+\mu v\\1\end{pmatrix}
=
\begin{pmatrix}p\\1\end{pmatrix}
+\mu
\begin{pmatrix}v\\0\end{pmatrix}.
\]
Cuando \(|\mu|\) crece, la dirección dominante es \((v,0)\). La última coordenada nula
identifica un punto ideal o punto en el infinito. Todas las rectas paralelas con dirección
\(v\) comparten ese mismo punto ideal.

\paragraph{Primera traducción algebraica.}
Una matriz proyectiva general produce
\[
\begin{pmatrix}X\\Y\\W\end{pmatrix}
=H
\begin{pmatrix}x\\y\\1\end{pmatrix}.
\]
Si \(W\neq0\), se vuelve al plano afín normalizando:
\[
  (x',y')=\left(\frac{X}{W},\frac{Y}{W}\right).
\]
Las afinidades son el caso especial en que la última fila de \(H\) es \((0,0,1)\), por lo
que \(W\) permanece igual a \(1\).

\begin{idea}[En palabras del profesor]
``Terminamos con dos formas de expresión algebraica, una en \(\R^2\) y otra en el plano
proyectivo. [\ldots] Sobre el espacio proyectivo nos preguntamos especialmente qué pasaba con
los puntos impropios y cómo se expresaría una transformación con puntos fijos'' (resumen de clase 3, cursada 2025).
\end{idea}

\subsection{La misma idea en cada dimensión}\label{sec:misma-idea}

\textbf{Desarrollo.} Casi todo lo que la materia hace en \(\R^2\), \(\R^3\) y \(\R^n\) ya
existe en \(\R\), donde es tan simple que ni se nota. Leer cada concepto primero en la
recta, después en el plano y después en el espacio es la versión más chica de ``ir a lo
chico para volver a lo grande'': si la idea se entiende con un número, la matriz sólo
agrega direcciones. La clase 10 del profesor hace exactamente esto con la inversa: ``la
inversa es un elemento de un conjunto que, al operar con otro, resulta en el elemento
neutro respecto de esa operación'', sea la suma, el producto, la composición o el producto
de matrices.

{\small
\begin{tabularx}{\textwidth}{>{\bfseries\raggedright\arraybackslash}p{17mm}>{\raggedright\arraybackslash}p{26mm}>{\raggedright\arraybackslash}p{32mm}>{\raggedright\arraybackslash}p{32mm}Y}
\toprule
& \(\R\) & \(\R^2\) & \(\R^3\) y \(\R^n\) & Hecho geométrico\\
\midrule
Identidad & el número \(1\): \(1\cdot a=a\) & la matriz \(I\): \(Iv=v\) & \(I\) de
\(n\times n\) & la transformación que no mueve nada; es el neutro de la composición\\
\addlinespace
Inversa & \(1/a\); no existe si \(a=0\) & \(A^{-1}\); no existe si \(\det A=0\) & igual;
si \(\det A=0\) hay núcleo & deshacer el movimiento; imposible si alguna dirección fue
aplastada a un punto (clase 10)\\
\addlinespace
Tamaño con signo & \(|a|\) estira; el signo dice si da vuelta & \(\det A\): área orientada
del paralelogramo imagen del cuadrado unidad & volumen orientado; en \(\R^n\),
hipervolumen & cuánto se agranda una figura y si cambia de orientación\\
\addlinespace
Producto & \(a\cdot b\) & producto interno \(u\cdot v=|u||v|\cos\theta\) & igual; en
homogéneas, \(\ell\cdot p=0\) es incidencia & proyección de un vector sobre otro; cero si
son perpendiculares\\
\addlinespace
Reflexión y rotación & \(x\mapsto -x\): la única ``rotación'' de la recta es media vuelta
& multiplicar por \(e^{i\theta}\) en \(\mathbb{C}=\R^2\); reflejar es conjugar &
matriz de rotación \(3\times3\) con un eje; reflexión en un plano & conservar distancias;
la rotación conserva la orientación, la reflexión la invierte\\
\addlinespace
Autovector & \(a\,x=\lambda x\): trivial, todo \(x\) sirve con \(\lambda=a\) &
\(Av=\lambda v\): sólo algunas direcciones, o ninguna real si hay giro & una rotación de
\(\R^3\) tiene exactamente una: el eje & dirección que cae sobre sí misma, sólo estirada o
invertida\\
\addlinespace
Derivada & \(f'(x)\): un número, la pendiente & gradiente: una forma lineal con dos
números; jacobiano \(2\times2\) si la función es vectorial & gradiente con \(n\) números;
jacobiano \(m\times n\) & la mejor aproximación lineal: recta tangente, plano tangente\\
\addlinespace
Recta & \(a\,x=b\): un punto & forma lineal \(n\cdot x=c\): una recta; al variar \(c\),
rectas paralelas & \(n\cdot x=c\): un plano; en \(\R^n\), hiperplano & el conjunto de
nivel de una regla graduada (sección~\ref{sec:formas-lineales})\\
\bottomrule
\end{tabularx}
}

\begin{idea}[Cómo usar la tabla]
Ante una fórmula nueva con matrices, escribirla primero con un solo número. Si en \(\R\)
la fórmula es obvia, la versión matricial dice lo mismo dirección por dirección; lo único
que agregan las dimensiones es que ahora las direcciones pueden mezclarse, y por eso
aparecen el orden de composición, los autovectores y el determinante.
\end{idea}

\section{Construcciones, cuentas y ejemplos}

Ahora construimos las afinidades geométricamente, las escribimos como \(F(x)=Ax+t\), buscamos sus puntos fijos y las seguimos sobre una figura concreta en la clase práctica. Los ejercicios de la guía son material docente; las cuentas y las figuras son desarrollo.

{\small\vspace{1mm}\linkUnidad{04_afines_homogeneas}{Simulación interactiva: Unidad 6 — Afinidades y coordenadas homogéneas}\par}

\subsection{Afinidades: idea geométrica y forma algebraica}\label{sec:afinidades-geom}

Una afinidad es una transformación que conserva el paralelismo. Como consecuencia, lleva
rectas en rectas y paralelogramos en paralelogramos. En coordenadas cartesianas adopta la forma
\[
  \boxed{F(x)=Ax+t},
\]
donde \(A\) es una matriz \(2\times2\) y \(t\in\R^2\). Para que sea invertible debe cumplirse
\(\det A\neq0\).

\begin{center}
\begin{tikzpicture}[scale=.86]
  \coordinate (a) at (1,1);
  \coordinate (b) at (4.1,1.2);
  \coordinate (c) at (5.2,3.2);
  \coordinate (d) at (2.1,3);
  \coordinate (ap) at (7,1.2);
  \coordinate (bp) at (10.4,.7);
  \coordinate (cp) at (11.7,3.4);
  \coordinate (dp) at (8.3,3.9);
  \draw[teal,line width=1.05pt] (a)--(b)--(c)--(d)--cycle;
  \draw[rust,line width=1.05pt] (ap)--(bp)--(cp)--(dp)--cycle;
  \draw[-{Latex[length=2.8mm]},tealdark,line width=.9pt]
    (5.7,2.2)--(6.6,2.2) node[midway,above] {$F$};
  \node[font=\small,text=teal] at (3.1,.55) {paralelogramo};
  \node[font=\small,text=rust] at (9.5,.25) {paralelogramo};
\end{tikzpicture}
\end{center}

\paragraph{Por qué conserva paralelas.}
Si dos rectas tienen la misma dirección \(v\), pueden escribirse
\[
  r_1(s)=p+sv,\qquad r_2(s)=q+sv.
\]
Sus imágenes son
\[
  F(r_1(s))=Ap+t+sAv,\qquad
  F(r_2(s))=Aq+t+sAv.
\]
Ambas comparten la dirección \(Av\); por lo tanto siguen siendo paralelas.

\subsection{Puntos fijos y forma afín}\label{sec:puntos-fijos-afin}

Esta sección enlaza la guía con el comienzo de la clase siguiente. Toda transformación afín
del plano puede escribirse
\[
  F(x)=Ax+t.
\]
Un punto fijo \(x_\ast\) satisface
\[
  F(x_\ast)=x_\ast
  \quad\Longleftrightarrow\quad
  (I-A)x_\ast=t.
\]

\begin{center}
\begin{tabularx}{\textwidth}{>{\bfseries}l l X}
\toprule
Transformación & \(A,t\) & Puntos fijos\\
\midrule
Traslación \(x+v\) & \(A=I,\ t=v\) & ninguno si \(v\neq0\)\\
Rotación de centro \(c\) & \(A=R,\ t=(I-R)c\) & \(c\), si el giro no es nulo\\
Homotecia de centro \(o\) & \(A=kI,\ t=(1-k)o\) & \(o\), si \(k\neq1\)\\
Reflexión en \(n\cdot x=c\) &
\(A=I-\dfrac{2nn^T}{\lVert n\rVert^2},\
t=\dfrac{2c}{\lVert n\rVert^2}n\) &
todos los puntos del eje; \(t\) es normal al eje\\
\bottomrule
\end{tabularx}
\end{center}

Si \(x_\ast\) existe, al trasladar el origen a ese punto mediante \(u=x-x_\ast\), la afinidad
queda puramente lineal:
\[
  F(x)-x_\ast=A(x-x_\ast)
  \qquad\Longrightarrow\qquad
  u'\!=Au.
\]
Esta es la idea detrás de buscar un sistema de coordenadas adecuado.

\subsection{Coordenadas homogéneas}\label{sec:homogeneas}

Las coordenadas homogéneas incorporan la traslación dentro de una multiplicación matricial.
El punto cartesiano \((x,y)\) se representa por \((x,y,1)\), y
\[
\begin{pmatrix}
x'\\y'\\1
\end{pmatrix}
=
\begin{pmatrix}
a&b&e\\
c&d&f\\
0&0&1
\end{pmatrix}
\begin{pmatrix}
x\\y\\1
\end{pmatrix}.
\]
La esquina superior izquierda contiene \(A\), mientras que la última columna contiene
el vector de traslación \(t=(e,f)^T\).

\begin{resultado}[Composición]
Si \(F_1(x)=A_1x+t_1\) y \(F_2(x)=A_2x+t_2\), entonces
\[
  F_2F_1(x)=A_2A_1x+(A_2t_1+t_2).
\]
En coordenadas homogéneas, esta regla queda incorporada en el producto ordinario de matrices
\(3\times3\).
\end{resultado}

\paragraph{Interpretación de la tercera coordenada.}
Los puntos ordinarios se representan con última coordenada distinta de cero y se normalizan
a \(w=1\). Los vectores de dirección o puntos ideales se representan con \(w=0\). Así, las
rectas paralelas pueden interpretarse proyectivamente como rectas que comparten un punto en
el infinito.

\subsection{Afinidades construidas geométricamente}\label{sec:afinidades-construidas}

El material de la segunda clase propone una afinidad definida por un eje fijo, una dirección
de desplazamiento y un factor. Los puntos del eje permanecen fijos; los demás se desplazan
sobre rectas paralelas a la dirección elegida.

\begin{center}
\begin{tikzpicture}[scale=.9]
  \draw[rust,line width=1.1pt] (.5,1)--(10.5,1) node[right] {eje};
  \coordinate (p) at (2.2,1);
  \coordinate (q) at (6.8,1);
  \coordinate (x) at (4.2,4.2);
  \coordinate (xp) at (7.2,3.1);
  \fill[ink] (p) circle (2.2pt) node[below] {$p$};
  \fill[ink] (q) circle (2.2pt) node[below] {$q$};
  \fill[teal] (x) circle (2.5pt) node[above left] {$x$};
  \fill[rust] (xp) circle (2.5pt) node[above right] {$F(x)$};
  \draw[mainline,-{Latex[length=2.5mm]}] (x)--(xp)
    node[midway,above,sloped,font=\small] {dirección};
  \draw[auxline] (p)--(x);
  \draw[auxline] (p)--(xp);
  \draw[auxline] (q)--(x);
  \draw[auxline] (q)--(xp);
\end{tikzpicture}
\end{center}

\begin{idea}[Una afinidad queda determinada por un triángulo]
Tres puntos no alineados forman un sistema de referencia afín. Si se especifican las imágenes
de los tres vértices y el triángulo imagen tampoco es degenerado, queda determinada una única
afinidad invertible que lleva el triángulo original al triángulo imagen. La razón es que dos
lados independientes determinan la acción lineal \(A\), y un vértice determina la traslación
\(t\).
\end{idea}

\subsection{Clase práctica: rotación y escala de una figura}

Para disponer de un caso reproducible, se toma una figura en forma de T con vértices
\[
\begin{aligned}
&(-2,2),\ (2,2),\ (2,1),\ (0.5,1),\\
&(0.5,-2),\ (-0.5,-2),\ (-0.5,1),\ (-2,1).
\end{aligned}
\]
La primera transformación gira \(45^\circ\) alrededor de \(c=(3,1)\). La segunda reduce a
la mitad respecto de \(d=(-1,2)\).

\begin{center}
\begin{tikzpicture}[scale=.68]
  \begin{scope}
    \draw[axisline,-{Latex[length=1.8mm]}] (-2.5,0)--(3.6,0);
    \draw[axisline,-{Latex[length=1.8mm]}] (0,-2.5)--(0,2.6);
    \draw[teal,line width=1.3pt]
      (-2,2)--(2,2)--(2,1)--(.5,1)--(.5,-2)--(-.5,-2)--(-.5,1)--(-2,1)--cycle;
    \fill[teal] (2,2) circle (2.5pt) node[labelback,above] {$q$};
    \fill[ink] (3,1) circle (2.5pt) node[labelback,right] {$c$};
    \node[labelback,font=\small] at (.3,-3.0) {figura inicial};
  \end{scope}

  \draw[-{Latex[length=2.5mm]},tealdark,line width=.9pt] (4.0,0)--(5.1,0)
    node[midway,above,font=\small] {\(R_{45^\circ,c}\)};

  \begin{scope}[xshift=6.2cm]
    \draw[axisline,-{Latex[length=1.8mm]}] (-2.5,0)--(3.8,0);
    \draw[axisline,-{Latex[length=1.8mm]}] (0,-4.0)--(0,2.0);
    \draw[blue,line width=1.3pt]
      (-1.243,-1.828)--(1.586,1.000)--(2.293,.293)--(1.232,-.768)--
      (3.354,-2.889)--(2.646,-3.596)--(.525,-1.475)--(-.536,-2.536)--cycle;
    \fill[blue] (1.586,1.000) circle (2.5pt) node[labelback,above] {$q_1$};
    \fill[ink] (3,1) circle (2.5pt) node[labelback,right] {$c$};
    \node[labelback,font=\small] at (.5,-4.5) {después de rotar};
  \end{scope}

  \draw[-{Latex[length=2.5mm]},tealdark,line width=.9pt] (10.2,0)--(11.3,0)
    node[midway,above,font=\small] {\(S_{0.5,d}\)};

  \begin{scope}[xshift=12.4cm]
    \draw[axisline,-{Latex[length=1.8mm]}] (-2.5,0)--(2.0,0);
    \draw[axisline,-{Latex[length=1.8mm]}] (0,-2.6)--(0,2.5);
    \draw[rust,line width=1.3pt]
      (-1.121,.086)--(.293,1.500)--(.646,1.146)--(.116,.616)--
      (1.177,-.445)--(.823,-.798)--(-.237,.263)--(-.768,-.268)--cycle;
    \fill[rust] (.293,1.500) circle (2.5pt) node[labelback,above] {$q_2$};
    \fill[ink] (-1,2) circle (2.5pt) node[labelback,left] {$d$};
    \fill[gold] (.293,.293) circle (2.8pt) node[labelback,right] {$p_\ast$};
    \node[labelback,font=\small] at (-.2,-3.0) {resultado final};
  \end{scope}
\end{tikzpicture}
\end{center}

\paragraph{Lectura geométrica.}
El punto señalado sigue
\[
  q=(2,2)
  \longmapsto
  q_1\approx(1.586,1.000)
  \longmapsto
  q_2\approx(0.293,1.500).
\]
La primera etapa es un arco conceptual alrededor de \(c\). La segunda transcurre sobre la
recta que une \(q_1\) con \(d\), reduciendo esa distancia a la mitad.

\begin{resultado}[Punto fijo de la composición]
El punto
\[
  p_\ast\approx(0.2929,0.2929)
\]
no es fijo para cada transformación por separado. Es fijo para la composición completa:
después de rotarlo alrededor de \(c\) y escalarlo respecto de \(d\), vuelve a la misma
posición.
\end{resultado}

\begin{idea}[Actividad de la clase]
Cada estudiante debe elegir su propia figura, coordenadas y centros. La comprobación debe
incluir la figura inicial, la imagen después de la rotación, la imagen después de la escala
y, si existe, el punto fijo de la composición.
\end{idea}

\subsubsection{Elementos a realizar}

\begin{itemize}[label=\(\square\)]
  \item Definir una figura propia mediante una tabla de coordenadas.
  \item Elegir un centro de rotación y otro centro para la escala.
  \item Construir gráficamente la figura inicial y las dos imágenes sucesivas.
  \item Escribir las matrices homogéneas \(H_1\), \(H_2\) y su producto \(H_2H_1\).
  \item Transformar los vértices en Excel mediante \texttt{MMULT}.
  \item Hallar el punto fijo y comprobar que la matriz compuesta lo deja invariado.
\end{itemize}

\subsection{Resumen de resultados}

\begin{center}
\begin{tabularx}{\textwidth}{>{\bfseries}p{31mm} X X}
\toprule
Objeto & Fórmula & Resultado clave\\
\midrule
Reflexión puntual &
\(\rho_b(x)=2b-x\) &
involución con único punto fijo \(b\)\\
Dos reflexiones puntuales &
\(\rho_a\rho_b(x)=x+2(a-b)\) &
traslación\\
Tres reflexiones puntuales &
\(\rho_a\rho_b\rho_c=\rho_d,\ d=a-b+c\) &
\(a,b,c,d\) forman paralelogramo\\
Reflexión en recta &
\(\varphi_B(x)=x-2\frac{n\cdot x-c}{\|n\|^2}n\) &
el eje queda fijo\\
Ejes paralelos &
\(\varphi_A\varphi_B(x)=x+2(\alpha-\beta)n\) &
traslación perpendicular\\
Ejes secantes &
\(\varphi_A\varphi_B=R_{o,2\theta}\) &
rotación alrededor de la intersección\\
Traslación y rotación &
\(T_vQ_c\neq Q_cT_v\) en general &
mismo ángulo, centros distintos\\
Homotecia &
\(\eta_o^k(x)=o+k(x-o)\) &
distancias multiplicadas por \(|k|\)\\
Similaridad directa &
\(S=\varphi_A\varphi_B\eta_o^\lambda\) &
escala y rotación con centro fijo\\
\bottomrule
\end{tabularx}
\end{center}

\section{En profundidad}

Para cerrar leemos la afinidad por sus autovalores y proponemos ejercicios de consolidación. Es desarrollo nuestro y prepara el cambio de base del capítulo~\ref{cap:08}.

\subsection{Autovalores y geometría de una transformación}\label{sec:autovalores-geom}

Una dirección \(v\neq0\) es propia de \(A\) si
\[
  Av=\lambda v.
\]
La dirección se conserva y solo cambia su escala o su sentido. Para hallar autovalores:
\[
  \det(A-\lambda I)=0.
\]
Geométricamente:
\begin{itemize}
  \item \(\lambda=1\) identifica una dirección que no cambia;
  \item \(\lambda=-1\) identifica una dirección que invierte su sentido;
  \item \(|\lambda|>1\) indica expansión;
  \item \(0<|\lambda|<1\) indica contracción.
\end{itemize}

En una rotación de \(\R^3\), el eje proporciona un autovector real con autovalor \(1\). Esta
observación conecta el estudio de rotaciones de la primera clase con las preguntas matriciales
de la clase de afinidades.

\begin{idea}[Preguntas guía para una matriz afín]
\begin{itemize}
  \item ¿Es invertible? Revisar \(\det A\).
  \item ¿Conserva orientación? Revisar el signo de \(\det A\).
  \item ¿Tiene puntos fijos? Resolver \((I-A)x=t\).
  \item ¿Tiene direcciones invariantes? Buscar autovectores de \(A\).
  \item ¿Cómo cambia áreas? El factor es \(|\det A|\).
\end{itemize}
\end{idea}

\subsection{Ejercicios adicionales de consolidación}

\begin{enumerate}
  \item Sean \(a=(2,-1)\), \(b=(-1,2)\) y \(x=(4,3)\). Calcule
  \(\rho_a\rho_b(x)\) paso a paso y verifique que coincide con
  \(x+2(a-b)\). Dibuje los cinco puntos relevantes.

  \item Dados \(a=(1,0)\), \(b=(3,1)\) y \(c=(5,4)\), construya
  \(d=a-b+c\). Verifique con pendientes o vectores que \(abcd\) es un
  paralelogramo.

  \item Refleje \(x=(3,4)\) respecto de \(B:2x-y=1\) usando:
  \begin{enumerate}[label=\alph*)]
    \item la fórmula con el vector normal;
    \item la proyección perpendicular;
    \item una construcción gráfica.
  \end{enumerate}

  \item Componga las reflexiones respecto de \(B:y=0\) y \(A:y=3\).
  Repita invirtiendo el orden. Compare los vectores de traslación.

  \item Construya dos rectas por el origen que formen \(30^\circ\).
  Multiplique sus matrices de reflexión y compruebe que resulta una rotación
  de \(60^\circ\).

  \item Para \(R=R(90^\circ)\) y \(v=(1,2)\), calcule las expresiones
  \(T_vR\) y \(RT_v\). Halle el punto fijo de cada una resolviendo
  \((I-R)c=t\).

  \item Dibuje un triángulo y aplique una homotecia de centro exterior con
  \(k=\frac32\). Compruebe que las longitudes se multiplican por
  \(\frac32\) y que los lados correspondientes son paralelos.

  \item Elija una escala \(\lambda=2\) y una rotación de \(40^\circ\) con
  centro \(o=(1,1)\). Escriba la similaridad como
  \(S(x)=o+2R(40^\circ)(x-o)\) y proponga dos ejes de reflexión cuya
  composición produzca el giro.
\end{enumerate}

\subsubsection{Respuestas para autocorrección}

\begin{enumerate}
  \item
  \[
  \rho_b(4,3)=(-6,1),\qquad
  \rho_a(-6,1)=(10,-3).
  \]
  Como \(2(a-b)=(6,-6)\), también
  \[
    (4,3)+(6,-6)=(10,-3).
  \]

  \item
  \[
    d=a-b+c=(3,3).
  \]
  Se verifica
  \[
  b-a=c-d=(2,1),
  \qquad
  c-b=d-a=(2,3).
  \]

  \item Con \(n=(2,-1)\), \(\lVert n\rVert^2=5\) y
  \(n\cdot(3,4)-1=1\):
  \[
  p=\left(\frac{13}{5},\frac{21}{5}\right),
  \qquad
  \varphi_B(3,4)=
  \left(\frac{11}{5},\frac{22}{5}\right).
  \]

  \item
  \[
  \varphi_A\varphi_B(x,y)=(x,y+6),
  \qquad
  \varphi_B\varphi_A(x,y)=(x,y-6).
  \]
  El orden invierte el sentido de la traslación.

  \item Tomando \(B\) como el eje \(x\) y \(A\) a \(30^\circ\):
  \[
  H_AH_B=
  \begin{pmatrix}
  \frac12&-\frac{\sqrt3}{2}\\
  \frac{\sqrt3}{2}&\frac12
  \end{pmatrix}
  =R(60^\circ).
  \]

  \item Para \(R=R(90^\circ)\):
  \[
  T_vR(x)=Rx+(1,2),
  \qquad
  RT_v(x)=Rx+(-2,1).
  \]
  Sus centros son, respectivamente,
  \[
    c_1=\left(-\frac12,\frac32\right),
    \qquad
    c_2=\left(-\frac32,-\frac12\right).
  \]

  \item Una elección concreta es
  \[
  o=(4,4),\quad
  A=(0,0),\quad B=(2,0),\quad C=(0,2).
  \]
  Con \(k=\frac32\):
  \[
  A'=(-2,-2),\quad B'=(1,-2),\quad C'=(-2,1).
  \]
  Los lados pasan de longitud \(2\) a \(3\), y los lados correspondientes permanecen
  paralelos.

  \item Pueden elegirse dos ejes que pasen por \(o=(1,1)\): \(B\) con dirección
  \(0^\circ\) y \(A\) con dirección \(20^\circ\). La composición
  \(\varphi_A\varphi_B\) gira \(40^\circ\); al combinarla con la homotecia de razón \(2\)
  se obtiene
  \[
    S(x)=o+2R(40^\circ)(x-o).
  \]
\end{enumerate}

\begin{idea}[Chequeo final para cada ejercicio]
\begin{itemize}
  \item ¿La figura muestra puntos iniciales, intermedios y finales?
  \item ¿Se distingue el orden de composición?
  \item ¿La cuenta usa la misma convención que el dibujo?
  \item ¿Se identifican puntos fijos y cantidades conservadas?
  \item ¿El ejemplo confirma exactamente la conclusión del teorema?
\end{itemize}
\end{idea}
